\section{What TeachQuiz measures}\label{sec:audit}

\paperfig{fig_failures.png}{F3}{\textbf{Why TeachQuiz cannot see teaching on MMMC.} (i) The grader lost answers whose ids copied the prompt's
template. (ii) The student refuses instead of forgetting. (iii) It already knows the answers. Bottom: our fix for
each.}{failures3}

TeachQuiz~\cite{code2video} gives a student a system prompt that forbids knowledge of the topic and allows
\texttt{NULL} (``insufficient evidence''). Our local student (\texttt{gemma4:26b}) sees 3 frames per section, since a
local model cannot take an mp4; Code2Video's protocol (whole video to Gemini) is implemented but not run here.

\paragraph{Grading.} The answer template's example id is \texttt{q1}. On MMMC's TEST items, named \texttt{t1..t5},
the student copied the template's ids in \OrderRuns{} of \AllRuns{} runs (\OrderShare\%); matched by id only, every
answer of those runs was scored \texttt{NULL}. DEV items are named \texttt{q1..} and were unaffected, which by itself
manufactures a DEV gain that cannot appear on TEST. We match by order when no id matches and the counts agree (M1).

\paragraph{Refusal.} Even when graded correctly, the student refuses \NullNoVideo\% of items without the video and
\NullWithVideo\% with it. $S_1$ is zero because the student obeys. A ``gain'' is then any change in its willingness to
answer.

\paragraph{Prior.} Table~\ref{tab:priors} removes NULL and shows the student nothing. It answers most MMMC items
correctly; the probabilities are near 1 and survive a cyclic shift of the options, so this is knowledge, not a letter
preference. Behind the NULL answers in $S_1$ is a prior ceiling, not a floor.

\begin{table*}[t]\centering\small
\caption{What the student already knows. Code2Video's unlearning prompt, but the student must choose A--D (no NULL). Accuracy: share of questions right. Confidence: mean probability the student puts on the right answer. Fully known: topics where every question is right with no video. Chance is 25.}\label{tab:priors}
\begin{tabular}{llrrrrr}\toprule Material & Student & Topics & Questions & Accuracy (\%) & Confidence (\%) & Fully known\\
\midrule
\multicolumn{7}{c}{\textit{MMMC (real topics from 3Blue1Brown)}}\\
No video & gemma4:26b & 117 & 585 & 99.0 & 99.0 & 108 \\
No video & qwen2.5vl:7b & 117 & 585 & 97.4 & 95.6 & 99 \\
No video & qwen2.5vl:3b & 117 & 585 & 91.9 & 86.6 & 67 \\
\midrule
\multicolumn{7}{c}{\textit{NOVEL-12 (invented concepts)}}\\
No video & gemma4:26b & 12 & 96 & 41.1 & 41.4 & 0 \\
Definition card as text & gemma4:26b & 12 & 96 & 84.9 & 83.7 & 1 \\
\bottomrule\end{tabular}
\end{table*}


\section{D-VGAIN and the metric ladder}\label{sec:ladder}

\paperfig{fig_measure.png}{F4}{\textbf{How we measure.} The same student, without NULL, answers the quiz from four materials. We read its
confidence in each letter, average two option orders, and report the gain over what it knew with no video.}{measure}

\paragraph{D-VGAIN.} The student must choose (no NULL), and we measure the prior instead of asking the student to
forget it. Four materials: none ($\Sblind$), TeachQuiz's 3 frames per section ($S_{\text{frames}}$), the first and
last frame of every section ($\Sstatic$, a slide deck of start and end states) and the silent video sampled at up to 4
frames per second ($\Svisual$; at most 410 frames, about 70 tokens each, so the prompt fits a 32k context, which we
check, because Ollama truncates silently). With $K = 4$ options, $c = 1/K$ and $\Beff = \max(c, \Sblind)$:
\begin{align}
  \mathrm{NVG}(S) &= \frac{S - \Beff}{1 - \Beff}, \qquad \mathrm{DML} = \Svisual - \Sstatic,\\
  \text{D-VGAIN} &= \mathrm{NVG}(\Svisual) + \lambda \max(0, \mathrm{DML}),
\end{align}
with $\lambda = 0.5$. NVG is Hake's normalised gain~\cite{hake1998} with the prior as pre-test. It is undefined when
less than $0.05$ of room is left ($1 - \Beff < 0.05$), and we report it as undefined, never as 0: with continuous
scores a no-video probability of $0.99999$ leaves $10^{-5}$ of room, and dividing by it turned noise into ``gains'' of
$\pm 10$ on the pilot's controls. DML credits what animation adds over its end states.

\paragraph{One probe, every version.} The student answers all questions of one material in one reply (``\texttt{1: B}'',
one line each) at temperature 0, and we read the log-probabilities of A--D at each letter's position: the top letter
is the forced-choice answer, $p(\text{right})$ the continuous score. Asking for a single letter per call failed:
the model's first token was often not a letter. The reply is repeated (at most twice, with a stricter line) when a
letter is missing, which happened when the student wrote an explanation of why it could not answer. Each material is
probed under two cyclic option orders and mapped back.

\paragraph{The ladder.} Each version of the measure adds one change: M0 TeachQuiz as published; M1 answers matched by order; M2 forced
choice ($S_{\text{frames}} - \Sblind$); M3 normalised gain; M4 the dense silent video; M5 the motion control
(D-VGAIN); M6 continuous scoring; M7 the mean of two option orders; M8 only items the student does not know without the video
($p_{\text{none}}(\text{right}) < 0.5$). Every version is judged by the same tests on the same videos:
\emph{defined} (share of videos with a score), \emph{separation} (a topic's own base video against another topic's
video, paired over topics; a teaching metric must put this above 0), \emph{reliability} (agreement of two measurements
over videos) and \emph{resolution} (ties among learner $-$ base differences).

\begin{table*}[t]\centering\footnotesize
\caption{Metric ladder on the MMMC pilot videos. 12 topics, 48 videos. Each row adds one change to the row above; points on a 0--100 scale. Scored: videos that get a score (-- : the student already knew everything). Own $-$ other: a topic's video minus another topic's video on the same quiz. Agreement: rank correlation of two measurements of the same videos. Learner $-$ \ldots: mean change over topics. Ties: topics where learner and base score the same.}\label{tab:ladder_mmmc}
\setlength{\tabcolsep}{4pt}\begin{tabular}{@{}llrrrrrrr@{}}\toprule & & \multicolumn{3}{c}{Does it measure teaching?} & \multicolumn{4}{c}{Does the learner repair help?}\\
\cmidrule(lr){3-5}\cmidrule(lr){6-9}
& Measure & Scored & \hd{Own $-$\\other ($\uparrow$)} & \hd{Agree-\\ment ($\uparrow$)} & \hd{Learner\\$-$ base} & \hd{Learner\\$-$ generic} & \hd{Learner $-$\\every section} & Ties ($\downarrow$)\\\midrule
M0 & TeachQuiz (Code2Video) & 48/48 & +10.6 & -- & +1.7 & $-$2.2 & $-$4.5 & 7/12 \\
M1 & + answers matched in order & 47/48 & +10.6 & 0.96 & +7.2 & +2.8 & +0.6 & 7/12 \\
M2 & + must choose (no NULL) & 48/48 & +16.7 & 1.00 & 0.0 & 0.0 & 0.0 & 12/12 \\
M3 & + gain over what it knew & 0/48 & -- & -- & -- & -- & -- & -- \\
M4 & + dense silent video & 0/48 & -- & -- & -- & -- & -- & -- \\
M5 & + motion bonus (D-VGAIN) & 0/48 & -- & -- & -- & -- & -- & -- \\
M6 & + confidence, not right/wrong & 0/48 & -- & -- & -- & -- & -- & -- \\
M7 & + two option orders & 0/48 & -- & -- & -- & -- & -- & -- \\
M8 & + only questions it did not know & 0/48 & -- & -- & -- & -- & -- & -- \\
M9 & M7 on 3 frames per section & 0/48 & -- & -- & -- & -- & -- & -- \\
\midrule
\multicolumn{9}{c}{\textit{Raw difference in confidence, not normalised (for reference)}}\\
R7 & Dense video minus no video & 48/48 & +4.1 & 0.65 & +1.4 & $-$0.1 & +0.3 & 1/12 \\
R9 & 3 frames minus no video & 48/48 & +13.0 & 0.77 & +0.1 & 0.0 & +0.1 & 0/12 \\
\bottomrule\end{tabular}
\end{table*}

